% !TeX root = ../教材题库_全部条目_转写工作稿.tex
% Source: XJ-B1-C05-TH-P192-Y1420
% 第5章 / 5.4 函数 y=A sin(ωx+φ) 的图象与性质 / 教材 p.192、193 / PDF 第199、200页
% OCR未校对：公式与文字均可能有误；未完成数学TeX转写；禁止直接发给学生
% 原题及相邻原解截图（核对用）：../题目原图/XJ-B1-C05-TH-P192-Y1420-01.webp, ../题目原图/XJ-B1-C05-TH-P192-Y1420-02.webp
% 本文件只在转写完成后作为题干使用；OCR 条目尚不是可靠题干。
\begin{bankitem}{p.192 思考（含紧邻说明）}{第5章 / 5.4 函数 y=A sin(\ensuremath{\omega}x+\ensuremath{\varphi}) 的图象与性质 / 教材 p.192、193 / PDF 第199、200页}{XJ-B1-C05-TH-P192-Y1420}{OCR未校对：公式与文字均可能有误；未完成数学TeX转写；禁止直接发给学生}
\begin{ocrdraft}
[教材p.192 / PDF第199页]
\par
？思考从函数3-工的图象出发，如何通过图象变换得到函数y=Asin(axr十@
\par
(A>O，>0）的图象？
\par
一般地，设A>O，>0，乒0，函数
\par
y=Asin(o\ocrglyph{00BB}x十@
\par
的图象可经过以下步骤得到：
\par
[教材p.193 / PDF第200页]
\par
（1）将正弦曲线y=sinx向左（当>0）或向右（当<0）平移鬲个单位长度；
\par
（2）将所得曲线上每一点的横坐标伸长（0<<1）或缩短\ensuremath{\text{\textperthousand}}>1）为原来的一（纵
\par
坐标不变）；
\par
（3）将所得曲线上每一点的纵坐标扩大（，A>1）或缩小（0<，A<1）为原来的A倍
\par
（横坐标不变）．
\par
函数十@的值域为[一，，A)，周期为
\end{ocrdraft}
\end{bankitem}
